\section{State of the art}
\label{sec:background}

\subsection{Swarm robotics}

Swarm robotics designs the collective behaviour of many simple robots that
act on local sensing and communication, without central control
\citep{brambilla2013}. Brambilla et al.\ classify the behaviours studied into
spatial organisation (aggregation, patterns), navigation (collective
exploration, coordinated motion, collective transport) and collective
decision making, and the design methods into behaviour-based design, by hand,
and automatic design. This project is a case of coordinated motion, flocking,
that carries a collective decision, where to go, made by a few informed
members; its controllers are designed by hand and studied in simulation.

\subsection{Flocking models}

\citet{reynolds1987} produced flocks for computer animation from three local
rules, each an acceleration: separation from neighbours that are too close,
alignment with their velocity, and cohesion toward their centre. Boids see
perfectly and respond at once, and they remain the usual starting point for
flocking robots: \citet{hauert2011} flew Reynolds' rules with a migration
point on ten fixed-wing robots, noting that a rotorcraft, which can move in
any direction, is the easiest robot to fly them on.

\citet{couzin2002} cast the same ingredients as zones around an individual: it
moves away from neighbours in the zone of repulsion, and only if there are
none, aligns with those in the zone of orientation and moves toward those in
the zone of attraction. Changing the widths of the zones makes a group switch
between a swarm, a milling torus and a parallel group, and the group's
history decides which, a collective memory.

\citet{olfatisaber2006} gave flocking a theory for agents with inertia: the
free rules generically break a flock into pieces, and a navigational
feedback, a virtual leader, holds it together and steers it; obstacles enter
as virtual agents.

Bringing the models to real flying robots meant modelling the robot.
\citet{viragh2014} describe a flying robot by its imperfections, the time it
takes to respond, its largest acceleration, delays, noise and limited range,
and show that delays make an ideal flocking model oscillate and collide; they
replace the stiff separation with a linear repulsion and add a pairwise
\emph{friction}, an alignment that damps velocity differences between close
neighbours. \citet{vasarhelyi2018} flew 30 drones with such a model in a
confined arena, its parameters optimised by evolution. Walls and obstacles
are virtual agents at their nearest point, flying away from them, with which
a drone aligns along a braking curve: the closer the wall, the less it may
fly toward it. Their arena has no goal: the flock only has to stay together
and inside.

Reactive models have limits in clutter. \citet{soria2021} show that the
potential-field family of swarm models does not guarantee fast and safe
flight among obstacles, and replace it with model predictive control.

\subsection{Leadership in groups on the move}

\citet{couzin2005} asked how animal groups move toward a goal that only a few
of their members know, when the others cannot tell who those are and nobody
signals. Their informed individuals are ordinary members of a zone-model
group, with one difference: they mix a preferred direction into the direction
their neighbours suggest, with a weight $\omega$. Three results frame this
project:
\begin{enumerate}
  \item a small informed minority can guide a large group, and the accuracy of
        the group grows with the fraction of informed individuals;
  \item the larger the group, the smaller the fraction needed for the same
        accuracy;
  \item with two informed subgroups preferring different directions, the
        group follows their average when the directions differ little, and
        chooses one of them, usually the larger subgroup's, when they differ
        a lot.
\end{enumerate}
The followers in this project are uninformed in the same sense: they don't
know where the goal is and can't recognise the leaders.

\subsection{Deliberative and hybrid control}

A purely reactive controller maps what a robot senses to what it does; a
deliberative one plans on a model of the world. Hybrid architectures put a
deliberative layer on top of a reactive one, which keeps the robot safe while
the plan gives it direction \citep{arkin1998}. Path finding on a grid with A*
\citep{hart1968}, with a cost that grows near obstacles, and path following by
pure pursuit, which steers toward a point a fixed distance ahead along the
path \citep{coulter1992}, are standard ways to fill the two layers in mobile
robotics. This project gives such a layer to the leaders only, on top of the
flocking model, and leaves the followers purely reactive.
